% QUESTAO
Sabemos que ao multiplicarmos uma matriz $A$, de ordem $n$, pela sua inversa $A^{-1}$, obtemos a matriz identidade $I_n$ (elemento neutro da multiplicação matricial), ou seja
\[A\cdot A^{-1}=I_n\textnormal{\quad e \quad}A^{-1}\cdot A=I_n \]
Partindo dessa definição, aplique o {\it teorema de Binet} para calcular $\det A^{-1}$, sabendo que $\det A=3$.

% ALTERNATIVAS
$\det A^{-1}=\frac{1}{3}$
$\det A^{-1}=\frac{2}{5}$
$\det A^{-1}=0$
$\det A^{-1}=1$
$\det A^{-1}=-3$


